\chapter{Conclusão}

O avanço acelerado da Inteligência Artificial, da Ciência de Dados, da Internet das Coisas (IoT), dos Grandes Modelos de Linguagem (LLMs), dos Sistemas Multiagentes e das arquiteturas computacionais distribuídas vem transformando profundamente a forma como organizações desenvolvem produtos, automatizam processos, tomam decisões e geram inovação. Nesse cenário, a formação de profissionais capazes de compreender, desenvolver, integrar e implantar soluções baseadas nessas tecnologias torna-se um fator estratégico para o desenvolvimento econômico, industrial e tecnológico.

O presente Projeto Pedagógico foi estruturado considerando as principais tendências tecnológicas nacionais e internacionais, incorporando conteúdos atualizados de Inteligência Artificial, Machine Learning, Deep Learning, Engenharia de Dados, Engenharia de Prompt, Engenharia de Contexto, IA Generativa, Sistemas Baseados em Grandes Modelos de Linguagem (LLMs), Agentes Inteligentes, Model Context Protocol (MCP), LangGraph, Sistemas Multiagentes, Visão Computacional, Internet das Coisas (IoT), Internet Industrial das Coisas (IIoT), Computação em Nuvem, MLOps, LLMOps, Governança de Inteligência Artificial e Segurança da Informação.

A organização curricular foi concebida de forma progressiva, permitindo que o estudante desenvolva competências desde os fundamentos científicos e computacionais até o desenvolvimento de soluções completas de Inteligência Artificial, integrando software, infraestrutura computacional, engenharia de dados, arquiteturas distribuídas e sistemas inteligentes autônomos. Essa abordagem favorece o desenvolvimento de profissionais preparados para atuar em diferentes segmentos econômicos, como indústria, tecnologia da informação, agronegócio, saúde, logística, comércio, serviços, energia e cidades inteligentes.

O curso também contempla competências relacionadas à ética, à governança de dados, à privacidade, à segurança da informação, à transparência algorítmica e ao uso responsável da Inteligência Artificial, aspectos indispensáveis para o desenvolvimento de soluções tecnológicas alinhadas às exigências regulatórias e às boas práticas internacionais.

A adoção de metodologias ativas, projetos integradores, laboratórios especializados, ambientes de desenvolvimento em nuvem e ferramentas amplamente utilizadas pelo mercado proporciona ao estudante uma formação prática, conectada aos desafios reais enfrentados pelas organizações. Ao longo do curso, o aluno desenvolve competências para atuar em equipes multidisciplinares, projetar arquiteturas inteligentes, integrar sistemas heterogêneos, automatizar processos e entregar soluções robustas, escaláveis e sustentáveis.

Sob a perspectiva institucional, esta proposta amplia a capacidade do SENAI de responder às transformações tecnológicas da indústria e da economia digital, fortalecendo sua atuação como agente de formação profissional alinhado às demandas emergentes do setor produtivo. A incorporação de temas como IA Generativa, Agentes Inteligentes, Sistemas Multiagentes, MCP, LLMOps e arquiteturas modernas de Inteligência Artificial posiciona a formação em sintonia com as tecnologias que vêm sendo adotadas por empresas líderes em inovação.

No contexto de Santa Catarina, estado reconhecido por seu ecossistema industrial diversificado, por seus polos de tecnologia e por seu crescente ambiente de inovação, este curso contribui para ampliar a disponibilidade de profissionais qualificados capazes de apoiar a transformação digital das empresas, fortalecer a competitividade da indústria catarinense e impulsionar o desenvolvimento de soluções tecnológicas de alto valor agregado.

Ao formar profissionais aptos a desenvolver aplicações inteligentes, sistemas autônomos e soluções integradas para diferentes setores econômicos, esta formação contribui para consolidar o ecossistema de inovação do estado, fortalecendo sua capacidade de adaptação às novas demandas da economia baseada em dados e Inteligência Artificial.

Dessa forma, este Projeto Pedagógico representa uma proposta educacional alinhada às transformações tecnológicas contemporâneas e às perspectivas futuras da computação inteligente, preparando profissionais para atuar em uma área estratégica para o desenvolvimento científico, tecnológico e industrial. Sua implementação reforça o compromisso institucional com a formação de excelência e contribui para manter Santa Catarina entre os estados brasileiros com maior capacidade de inovação, desenvolvimento tecnológico e qualificação profissional em Inteligência Artificial.
